% GENERATED FILE. Edit canonical lessons, then run npm run generate:books.
% canonical-source-hash: fnv1a64:43ae4a2f8657a613
% canonical-lessons: HI-C308-bera, HI-C308-gotakhor, HI-C308-zinda, HI-C308-churana, HI-C308-lapakna

\chapter{The Rescue}
\label{ch:hi-the-rescue}
\hlchaptermodality{308}

\begin{chapteropening}
\textbf{By the end of this chapter:} \emph{I can tell how people were rescued.}

Complete the last lesson of chapter 308: I can tell how people were rescued.
\end{chapteropening}

\section[\texorpdfstring{\hi{बेड़ा} (beṛā)}{बेड़ा (beṛā)}]{\hi{बेड़ा} (beṛā) — a raft; a fleet}

\label{lesson:HI-C308-bera}



\begin{quote}
Before the new one: say the Hindi for to take out, then the Hindi for in a sorry state.
\end{quote}

\subsection*{You'll want to know: \hi{बेड़ा}}
\textbf{\hi{बेड़ा}} — \emph{beṛā} — \textquotedblleft{}a raft; a fleet\textquotedblright{}.

\begin{quote}
\hi{लोगों} \hi{ने} \hi{लकड़ी} \hi{का} \hi{एक} \hi{बेड़ा} \hi{बनाया} \hi{और} \hi{बच्चों} \hi{को} \hi{पानी} \hi{से} \hi{निकाला।}
\end{quote}

\begin{quote}
\emph{logoṁ ne lakṛī kā ek beṛā banāyā aur baccoṁ ko pānī se nikālā.}
\end{quote}

\begin{quote}
People built a wooden raft and got the children out of the water.
\end{quote}

Masculine.

\begin{cousinweb}
goes back to Sanskrit \emph{bēḍā}, \textquotedblleft{}a boat\textquotedblright{} (CDIAL 9308), where CDIAL gives Hindi \emph{beṛā}, a boat or a raft. In a flood a \textbf{\hi{बेड़ा}} is anything that floats and carries people; with ships, it is a fleet.
\end{cousinweb}

\begin{grammarlens}[title={What you've built}]
One more word for telling what happened.
\end{grammarlens}

\subsection*{Your turn}
\begin{itemize}
\raggedright
  \item \emph{Say it:} \emph{beṛā}
  \item \emph{Say it:} \emph{beṛā}, once more
  \item \emph{Say it:} \emph{logoṁ ne lakṛī kā ek beṛā banāyā aur baccoṁ ko pānī se nikālā.} — People built a wooden raft and got the children out of the water.
  \item \emph{Recall:} say the Hindi for to take out, then the Hindi for in a sorry state, then say \emph{beṛā} again
\end{itemize}

\begin{tcolorbox}[breakable,colback=teal!4,colframe=teal!35!black,title={Before you move on}]
What does \hi{बेड़ा} mean? (\textquotedblleft{}A raft; a fleet\textquotedblright{}.) Say it once more.
\end{tcolorbox}
\section[\texorpdfstring{\hi{गोताख़ोर} (gotākhor)}{गोताख़ोर (gotākhor)}]{\hi{गोताख़ोर} (gotākhor) — a diver}

\label{lesson:HI-C308-gotakhor}



\begin{quote}
Before the new one: say the Hindi for in a sorry state, then the Hindi for a raft.
\end{quote}

\subsection*{You'll want to know: \hi{गोताख़ोर}}
\textbf{\hi{गोताख़ोर}} — \emph{gotākhor} — \textquotedblleft{}a diver\textquotedblright{}.

\begin{quote}
\hi{गोताख़ोर} \hi{कितने} \hi{बजे} \hi{पहुँचेंगे}?
\end{quote}

\begin{quote}
\emph{gotākhor kitne baje pahũceṅge?}
\end{quote}

\begin{quote}
What time will the divers get here?
\end{quote}

Masculine or feminine, by the person meant.

\begin{cousinweb}
\textbf{\hi{गोताख़ोर}} (\emph{gotākhor}) is usually traced to Persian \emph{ġōṭa-xōr}, built from \emph{ġōṭa}, a plunge, and \emph{-xōr}, one who takes: someone who dives under the water.
\end{cousinweb}

\begin{grammarlens}[title={What you've built}]
One more word for telling what happened.
\end{grammarlens}

\subsection*{Your turn}
\begin{itemize}
\raggedright
  \item \emph{Say it:} \emph{gotākhor}
  \item \emph{Say it:} \emph{gotākhor}, once more
  \item \emph{Say it:} \emph{gotākhor kitne baje pahũceṅge?} — What time will the divers get here?
  \item \emph{Recall:} say the Hindi for in a sorry state, then the Hindi for a raft, then say \emph{gotākhor} again
\end{itemize}

\begin{tcolorbox}[breakable,colback=teal!4,colframe=teal!35!black,title={Before you move on}]
What does \hi{गोताख़ोर} mean? (\textquotedblleft{}A diver\textquotedblright{}.) Say it once more.
\end{tcolorbox}
\section[\texorpdfstring{\hi{ज़िंदा} (zindā)}{ज़िंदा (zindā)}]{\hi{ज़िंदा} (zindā) — alive}

\label{lesson:HI-C308-zinda}



\begin{quote}
Before the new one: say the Hindi for a raft, then the Hindi for a diver.
\end{quote}

\subsection*{You'll want to know: \hi{ज़िंदा}}
\textbf{\hi{ज़िंदा}} — \emph{zindā} — \textquotedblleft{}alive\textquotedblright{}.

\begin{quote}
\hi{गोताख़ोरों} \hi{ने} \hi{नदी} \hi{से} \hi{एक} \hi{आदमी} \hi{को} \hi{ज़िंदा} \hi{निकाला।}
\end{quote}

\begin{quote}
\emph{gotākhoroṁ ne nadī se ek ādmī ko zindā nikālā.}
\end{quote}

\begin{quote}
The divers pulled a man out of the river alive.
\end{quote}

\begin{cousinweb}
\textbf{\hi{ज़िंदा}} (\emph{zindā}) is usually traced to Persian \emph{zinda}, \textquotedblleft{}alive, living\textquotedblright{}. It does not change for a woman or for many: \textbf{\hi{वह} \hi{ज़िंदा} \hi{है}}, \textbf{\hi{वे} \hi{ज़िंदा} \hi{हैं}}.
\end{cousinweb}

\begin{grammarlens}[title={What you've built}]
One more word for telling what happened.
\end{grammarlens}

\subsection*{Your turn}
\begin{itemize}
\raggedright
  \item \emph{Say it:} \emph{zindā}
  \item \emph{Say it:} \emph{zindā}, once more
  \item \emph{Say it:} \emph{gotākhoroṁ ne nadī se ek ādmī ko zindā nikālā.} — The divers pulled a man out of the river alive.
  \item \emph{Recall:} say the Hindi for a raft, then the Hindi for a diver, then say \emph{zindā} again
\end{itemize}

\begin{tcolorbox}[breakable,colback=teal!4,colframe=teal!35!black,title={Before you move on}]
What does \hi{ज़िंदा} mean? (\textquotedblleft{}Alive\textquotedblright{}.) Say it once more.
\end{tcolorbox}
\section[\texorpdfstring{\hi{छुड़ाना} (chuṛānā)}{छुड़ाना (chuṛānā)}]{\hi{छुड़ाना} (chuṛānā) — to set free; to get loose}

\label{lesson:HI-C308-churana}



\begin{quote}
Before the new one: say the Hindi for a diver, then the Hindi for alive.
\end{quote}

\subsection*{You'll want to know: \hi{छुड़ाना}}
\textbf{\hi{छुड़ाना}} — \emph{chuṛānā} — \textquotedblleft{}to set free; to get loose\textquotedblright{}.

\begin{quote}
\hi{अरुण} \hi{ने} \hi{रस्सी} \hi{काटकर} \hi{बकरी} \hi{को} \hi{छुड़ाया।}
\end{quote}

\begin{quote}
\emph{Aruṇ ne rassī kāṭkar bakrī ko chuṛāyā.}
\end{quote}

\begin{quote}
Arun cut the rope and set the goat free.
\end{quote}

Past, with \hi{ने} (he or she …): \hi{उसने} \hi{बकरी} \hi{को} \hi{छुड़ाया}, \emph{usne bakrī ko chuṛāyā}.

\begin{cousinweb}
\textbf{\hi{छुड़ाना}} (\emph{chuṛānā}) is the causative of \textbf{\hi{छूटना}}, to be released, which CDIAL 3707 sets under \emph{*kṣuṭyatē}, \textquotedblleft{}is released\textquotedblright{}: to make something come loose.
\end{cousinweb}

\begin{grammarlens}[title={What you've built}]
One more word for telling what happened.
\end{grammarlens}

\subsection*{Your turn}
\begin{itemize}
\raggedright
  \item \emph{Say it:} \emph{chuṛānā}
  \item \emph{Say it:} \emph{chuṛānā}, once more
  \item \emph{Say it:} \emph{Aruṇ ne rassī kāṭkar bakrī ko chuṛāyā.} — Arun cut the rope and set the goat free.
  \item \emph{Recall:} say the Hindi for a diver, then the Hindi for alive, then say \emph{chuṛānā} again
\end{itemize}

\begin{tcolorbox}[breakable,colback=teal!4,colframe=teal!35!black,title={Before you move on}]
What does \hi{छुड़ाना} mean? (\textquotedblleft{}To set free; to get loose\textquotedblright{}.) Say it once more.
\end{tcolorbox}
\section[\texorpdfstring{\hi{लपकना} (lapaknā)}{लपकना (lapaknā)}]{\hi{लपकना} (lapaknā) — to dart; to spring forward}

\label{lesson:HI-C308-lapakna}



\begin{quote}
Before the new one: say the Hindi for alive, then the Hindi for to set free.
\end{quote}

\subsection*{You'll want to know: \hi{लपकना}}
\textbf{\hi{लपकना}} — \emph{lapaknā} — \textquotedblleft{}to dart; to spring forward\textquotedblright{}.

\begin{quote}
\hi{मीरा} \hi{ने} \hi{लपककर} \hi{बच्चे} \hi{को} \hi{पकड़ा।}
\end{quote}

\begin{quote}
\emph{Mīrā ne lapakkar bacce ko pakṛā.}
\end{quote}

\begin{quote}
Mira darted forward and caught the child.
\end{quote}

Past (he …): \hi{वह} \hi{लपका}, \emph{vah lapkā}.

\begin{cousinweb}
\textbf{\hi{लपकना}} (\emph{lapaknā}) is inherited from \emph{*lappa}, a form CDIAL reconstructs for leaping (CDIAL 10939). It is a quick spring towards something: \textbf{\hi{लपककर} \hi{पकड़ा}}, sprang and caught.
\end{cousinweb}

\begin{grammarlens}[title={What you've built}]
One more word for telling what happened.
\end{grammarlens}

\subsection*{Your turn}
\begin{itemize}
\raggedright
  \item \emph{Say it:} \emph{lapaknā}
  \item \emph{Say it:} \emph{lapaknā}, once more
  \item \emph{Say it:} \emph{Mīrā ne lapakkar bacce ko pakṛā.} — Mira darted forward and caught the child.
  \item \emph{Recall:} say the Hindi for alive, then the Hindi for to set free, then say \emph{lapaknā} again
\end{itemize}

\begin{tcolorbox}[breakable,colback=teal!4,colframe=teal!35!black,title={Before you move on}]
What does \hi{लपकना} mean? (\textquotedblleft{}To dart; to spring forward\textquotedblright{}.) Say it once more.
\end{tcolorbox}
